\documentclass[11pt]{article}
\usepackage[T1]{fontenc}
\usepackage{lmodern}
\usepackage{microtype}
\usepackage[a4paper,margin=2.3cm]{geometry}
\usepackage{amsmath,amssymb}
\usepackage{graphicx}
\usepackage{booktabs}
\usepackage{enumitem}
\usepackage{xcolor}
\usepackage{caption}
\usepackage[most]{tcolorbox}
\usepackage{fancyhdr}
\usepackage{titlesec}
\usepackage[colorlinks=true,linkcolor=azul,citecolor=azul,urlcolor=azul]{hyperref}
\definecolor{azul}{HTML}{1D4ED8}
\definecolor{tinta}{HTML}{14213D}
\definecolor{ambar}{HTML}{B45309}
\definecolor{gris}{HTML}{64748B}
\definecolor{fondo}{HTML}{EEF2F7}
\graphicspath{{../../02_RESULTADOS/Graficos/presentacion/}{../../02_RESULTADOS/Animaciones_MP4/}}
\renewcommand{\figurename}{Figura}
\renewcommand{\tablename}{Tabla}
\renewcommand{\contentsname}{Índice}
\renewcommand{\refname}{Referencias}
\renewcommand{\appendixname}{Anexo}
\titleformat{\section}{\Large\bfseries\color{tinta}}{\thesection}{0.8em}{}
\titleformat{\subsection}{\large\bfseries\color{tinta}}{\thesubsection}{0.8em}{}
\setlist{itemsep=2pt,topsep=4pt}
\captionsetup{font=small,labelfont=bf}
\newtcolorbox{idea}[1][]{colback=fondo,colframe=fondo,coltitle=tinta,fonttitle=\bfseries,title=#1,sharp corners,boxrule=0pt,left=6pt,right=6pt,top=4pt,bottom=4pt}
\newtcolorbox{alerta}[1][]{colback=ambar!8,colframe=ambar,fonttitle=\bfseries,title=#1,sharp corners,boxrule=0.6pt,left=6pt,right=6pt}
\newcommand{\md}{\dot m}
\newcommand{\Qgw}{\dot Q_{gw}}
\newcommand{\Qwa}{\dot Q_{wa}}
\newcommand{\cvi}{c_v^{ig}}
\newcommand{\dd}{\mathrm{d}}
\pagestyle{fancy}
\fancyhf{}
\renewcommand{\headrulewidth}{0pt}
\fancyfoot[L]{\small\color{gris}ICF130 · Tema 7 · material interno del grupo}
\fancyfoot[R]{\small\color{gris}\thepage}
\setlength{\parindent}{0pt}
\setlength{\parskip}{5pt}

\title{\color{tinta}\textbf{Guía de exposición}\\[4pt]\large Presentación final · ICF130 Tema 7}
\author{Material interno del grupo}
\date{Noviembre 2026}
\begin{document}
\maketitle
\section{Formato}
\begin{itemize}
\item 21 diapositivas (la 21 son referencias, no se expone). Duración objetivo: \textbf{18 a 20 minutos} más preguntas,
a 1 minuto por diapositiva aprox. Si piden acortar a 12-15 min, ver la sección~\ref{sec:corto}.
\item Las notas del presentador de cada diapositiva (en el .pptx, vista Moderador) tienen el texto oral sugerido.
No leerlas: usarlas como guion de ideas.
\item El video de la diapositiva 18 va incrustado en el .pptx (18 s). Si el proyector falla, abrir
\texttt{02\_RESULTADOS/Animaciones\_MP4/ICF130\_V1\_V2\_V3\_tres\_EOS.mp4} directamente. El PDF muestra solo el último cuadro.
\end{itemize}

\section{División entre integrantes}
\begin{center}\small
\begin{tabular}{llll}\toprule
Expositor & Diapositivas & Contenido & Tiempo\\\midrule
Paul & 1-5 & Portada, problema, hipótesis, balances, etapas & 5-6 min\\
Benjamín & 6-12 & EOS, energía, caudal, método, verificación, tabla, V3 & 7 min\\
Andrés & 13-20 & EOS reales, pared, isotermo, suministro, Span, video, cierre & 7 min\\\bottomrule
\end{tabular}\end{center}
La división pone la física de base en Paul, la matemática y la numérica en Benjamín (lo más técnico, con la
tabla de resultados como cierre de su parte) y la interpretación en Andrés.

\section{Diapositiva por diapositiva: lo que no puede faltar}
\begin{enumerate}[leftmargin=*]
\item \textbf{Portada (Paul).} Tema 7; un mismo modelo en tres etapas y con tres EOS.
\item \textbf{Problema.} Entra a 27 °C y llega a 146 °C sin calor: el trabajo de flujo. Tres preguntas.
\item \textbf{Modelo 0D.} Una sola $P$ y $T$; rígido, sin trabajo de eje; caso base. Los parámetros de pared son de estudio.
\item \textbf{Balances.} Las tres ecuaciones; $h=u+Pv$; solución exacta $\gamma T_s\approx420$ K usada para verificar.
\item \textbf{Etapas.} Un fenómeno por etapa; tres EOS en condiciones idénticas. \emph{Transición:} ``Benjamín explica cómo entran las ecuaciones de estado''.
\item \textbf{EOS (Benjamín).} Qué corrige cada una; mismas constantes críticas y mismo $\cvi$.
\item \textbf{Energía consistente.} Maxwell $\Rightarrow$ $u$ depende de la EOS; EDO de $T$; recupera $2am/V$ y el ideal.
\item \textbf{Caudal.} Estrangulado hasta 106 bar; misma ley para las tres EOS (hipótesis declarada).
\item \textbf{Método.} Brent en rama estable; RK45; evento; cuadratura independiente.
\item \textbf{Verificación.} $2{,}6\cdot10^{-6}$ K, orden 4,03, residuo $<10^{-11}$, segunda formulación. Verificación $\neq$ validación.
\item \textbf{Tabla.} V3: PR llega 22,8 s antes y 6,7 K más frío. En V1 más masa, en V2-V3 menos.
\item \textbf{V3.} Escalas 2 s y 360 s; pico del ideal; la pared almacena. \emph{Transición:} ``Andrés explica por qué difieren''.
\item \textbf{Por qué difieren (Andrés).} Menor $h_s$ (311,7 vs 275,4 kJ/kg) y $Z>1$ a 190 bar.
\item \textbf{Sensibilidad.} $UA_{gw}$: 0 = V2, 8 = base; más intercambio, más frío y más lento; misma tendencia en las tres EOS.
\item \textbf{Isotermo.} 940 kJ; 40 \% más masa y tiempo; V3 entre los extremos.
\item \textbf{Suministro.} El signo cambia; a 200 bar/300 K dos efectos de 26 s se cancelan.
\item \textbf{Span.} 5,68 / 4,90 / 2,83 \%; PR la más cercana, no exacta.
\item \textbf{Video.} Color = $T_g$ media; anillo = $T_w$; barra = $P/P_{obj}$; por qué se quitó el mapa anterior.
\item \textbf{Límites.} 0D, conexión ideal, parámetros no medidos, sin experimento.
\item \textbf{Conclusiones.} Cuatro ideas; cerrar con ``quedamos atentos a sus preguntas''.
\end{enumerate}

\section{Versión corta (12-15 min)}\label{sec:corto}
Ocultar las diapositivas 9 (método), 14 (sensibilidad) y 15 (isotermo) y resumirlas en una frase en la 10 y la 12.
Quedan 17 diapositivas expuestas.

\section{Antes de exponer}
\begin{itemize}
\item Abrir el .pptx en el computador de la sala y probar que el video se reproduce.
\item Llevar el PDF y el MP4 en un pendrive por si falla PowerPoint.
\item Tener a mano \texttt{Preguntas\_Profesor.pdf}: cada uno estudia las preguntas de su parte y las difíciles (27-37).
\item Si piden ver el código: \texttt{python main.py --caso V3 --eos pr} imprime un resumen en segundos; \texttt{python -m pytest} corre las 25 pruebas en unos 10 s.
\end{itemize}

\section{Cifras para memorizar}
\begin{center}\small
\begin{tabular}{ll}\toprule
Cambio de régimen & 105,66 bar ($r^*=0{,}528$)\\
V2 ideal & 418,92 s · 419,09 K (146 °C)\\
V3 ideal / VdW / PR & 456,40 / 440,75 / 433,60 s · 385,32 / 379,44 / 378,59 K\\
Masa V3 ideal / PR & 8,307 / 7,985 kg\\
$Z$ final V3 VdW / PR & 1,037 / 1,059\\
$h_s$ ideal / VdW / PR & 311,7 / 273,5 / 275,4 kJ/kg\\
Calor a la pared V3 ideal / PR & 208,6 / 161,8 kJ\\
Isotermo: tiempo y calor (ideal) & 586,4 s · 945 kJ\\
Error densidad vs Span (200 bar, 300 K) & 5,68 / 4,90 / 2,83 \%\\
Orden RK4 & 4,03\\\bottomrule
\end{tabular}\end{center}
\end{document}
